\documentclass[10pt]{amsart}
\usepackage[a4paper,margin=23mm]{geometry}
\usepackage{amsmath,amssymb,amsthm,xfrac,hyperref,xspace,graphicx}
\usepackage[english]{babel}
\usepackage[scr=rsfs,cal=euler]{mathalfa}
\newcommand{\resp}{resp.}
\newcommand{\cf}{cf.}
\newcommand{\ie}{i.e.}
\newcommand{\eg}{e.g.}
\newcommand{\loccit}{loc.\ cit.}
\newcommand{\skpt}{\par\smallskip\noindent}
\let\Subsection\subsection
\let\Subsubsection\subsubsection
\emergencystretch=3em
\pagestyle{plain}
\multlinegap0pt

\hyphenation{after approach char-ac-ter exam-ple involv-ing inspec-tion irre-du-cible}

\makeatletter
\def\l@part{\@tocline{-1}{5pt plus2pt}{0pt}{}{}}%
\def\@auxsection{\part}
\makeatother

\newenvironment{enumeratei}
{\bgroup\def\theenumi{\roman{enumi}}\def\theenumii{\arabic{enumii}}\begin{enumerate}}
{\end{enumerate}\egroup}
\newenvironment{enumeratea}
{\bgroup\def\theenumi{\alph{enumi}}\def\theenumii{\roman{enumii}}\begin{enumerate}}
{\end{enumerate}\egroup}

\newcommand{\spfrac}[2]{\sfrac{#1}{(#2)}}
\newcommand\from{\mathchoice{\longleftarrow}{\leftarrow}{\leftarrow}{\leftarrow}}
\newcommand\mto{\mathchoice{\longmapsto}{\mapsto}{\mapsto}{\mapsto}}
\newcommand\hto{\mathrel{\lhook\joinrel\to}}

\newcommand{\Changel}{\mathcode`l="7160}
\newcommand{\Changelback}{\mathcode`l="716C}
\newcommand{\NoChangel}[1]{%
\expandafter\let\csname old\string#1\endcsname=#1
\let#1=\relax
\newcommand{#1}{\mathop{\mathcode`l="716C\csname old\string#1\endcsname\mathcode`l="7160}\displaylimits}%
}

\NoChangel{\log}
\NoChangel{\ln}
\NoChangel{\lim}
\NoChangel{\limsup}
\NoChangel{\liminf}
\NoChangel{\varlimsup}
\NoChangel{\varliminf}
\NoChangel{\varinjlim}
\NoChangel{\varprojlim}


\usepackage{booktabs}
\usepackage{graphicx}
%\usepackage[table]{xcolor}
\usepackage{colortbl}
\usepackage{dynkin-diagrams}
\usetikzlibrary{babel}

\usepackage{amscd}

\newcommand{\proofsec}[1]{\subsubsection*{#1}}
\newcounter{stepnum}
\newcommand{\thestep}{\arabic{stepnum}}
\newcommand{\step}[1]{\refstepcounter{stepnum}\proofsec{Step \textup{\thestep}} \emph{#1}}

\renewcommand{\thepart}{\Roman{part}}

\newtheorem{thm}[equation]{Theorem}
\newtheorem{lem}[equation]{Lemma}
\newtheorem{cor}[equation]{Corollary}
\newtheorem{prop}[equation]{Proposition}
\newtheorem{conj}[equation]{Conjecture}

\theoremstyle{definition}
\newtheorem{defn}[equation]{Definition}
\newtheorem{example}[equation]{Example}
\newtheorem{egs}[equation]{Examples}
\newtheorem{rmk}[equation]{Remark}
\newtheorem{rmks}[equation]{Remarks}
\newtheorem*{rmk*}{Remark}


\let\eps\varepsilon
\let\grad\nabla
\let\ic\circ
\let\la\lambda
\let\laplacian\Delta
\let\och\varpi
\let\sinfty c
\let\vf\delta
\let\:\colon

\newcommand{\mfa}{\mathfrak{a}}
\newcommand{\mfe}{\mathfrak{e}}
\newcommand{\mfg}{\mathfrak{g}}
\newcommand{\mfm}{\mathfrak{m}}
\newcommand{\mft}{\mathfrak{t}}

\newcommand{\trans}{\mathsf{T}}

\newcommand{\A}{\mathbb{A}}
\newcommand{\R}{\mathbb{R}}
\newcommand{\F}{\mathbb{F}}
\newcommand{\Q}{\mathbb{Q}}
\newcommand{\E}{\mathbb{E}}
\newcommand{\C}{\mathbb{C}}
\newcommand{\Z}{\mathbb{Z}}
\newcommand{\Gm}{\mathbb{G}_{\mathrm{m}}}

\newcommand{\mcD}{\mathcal{D}}
\newcommand{\mcO}{\mathcal{O}}
\let\cO\mcO
\let\sO\cO
\newcommand{\cOh}{\hat{\mcO}}

\renewcommand{\ss}{\mathrm{ss}}

\newcommand{\hst}{\widetilde{\alpha}}

 \newcommand{\ditto}{\ \textquotedbl \ }
\newcommand{\slot}{\--}
\newcommand{\emptyslot}{{-}}
\newcommand{\qform}[1]{{\left\langle{#1}\right\rangle}}
\newcommand{\eand}{\quad\text{and}\quad}
\newcommand{\iso}{\overset{\sim}\to}
\newcommand{\mmin}{\underline{m}}
\newcommand{\mmax}{\overline{m}}

\renewcommand{\div}{\mathop{\mathrm{div}}}

\DeclareMathOperator{\RE}{Re}
\DeclareMathOperator{\IM}{Im}
\renewcommand{\Re}{\RE}
\DeclareMathOperator{\USp}{USp}
\DeclareMathOperator{\GL}{GL}
\DeclareMathOperator{\PGL}{PGL}
\DeclareMathOperator{\U}{U}
\DeclareMathOperator{\AGL}{AGL}
\DeclareMathOperator{\SL}{SL}
\DeclareMathOperator{\Spin}{Spin}
\DeclareMathOperator{\Sp}{Sp}
\DeclareMathOperator{\SU}{SU}
\DeclareMathOperator{\SO}{SO}
\DeclareMathOperator{\Spec}{Spec}
\DeclareMathOperator{\class}{Class}
\newcommand{\ClG}{\class(G)}
\DeclareMathOperator{\End}{End}
\DeclareMathOperator{\Tr}{Tr}
\DeclareMathOperator{\Ad}{Ad}
\DeclareMathOperator{\Lie}{Lie}
\DeclareMathOperator{\Aut}{Aut}
\DeclareMathOperator{\Int}{Int}
\DeclareMathOperator{\Out}{Out}
\DeclareMathOperator{\rank}{rank}
\DeclareMathOperator{\Lin}{Lin}
\DeclareMathOperator{\im}{im}
\DeclareMathOperator{\Sym}{Sym}
\DeclareMathOperator{\Hom}{Hom}
\DeclareMathOperator{\Res}{Res}
\DeclareMathOperator{\Der}{Der}
\DeclareMathOperator{\Reg}{Reg}
\DeclareMathOperator{\tr}{Tr}
\DeclareMathOperator{\sgn}{sgn}
\DeclareMathOperator{\diag}{diag}
\DeclareMathOperator{\Diff}{Diff}


\newcommand{\kh}{\hat{k}}
\newcommand{\gh}{\hat{g}}
\newcommand{\Gh}{\hat{G}}
\newcommand{\Gs}{\bar{G}}
\newcommand{\Mh}{\hat{M}}
\newcommand{\kalg}{\bar{k}}
\newcommand{\Qalg}{\bar{\Q}}

\newcommand{\Tt}{\tilde{T}}

\newcommand{\ov}{\omega^\vee}
\newcommand{\av}{\alpha^\vee}
\newcommand{\Qv}{Q^\vee}
\newcommand{\Pv}{P^\vee}
\newcommand{\Rv}{R^\vee}

\newcommand{\cartan}[1]{\mathsf{#1}}
\newcommand{\At}{\cartan{A}}
\newcommand{\Bt}{\cartan{B}}
\newcommand{\Ct}{\cartan{C}}
\newcommand{\Gt}{\cartan{G}}

\newcommand{\thinfty}{\theta_0}

\newcommand{\Wa}{W_a}

\def\edgelen{1cm} % for the Dynkin diagrams

\numberwithin{equation}{section}
\begin{document}
\title{Uniform negative normalized character values for nontrivial irreducible representations of each connected compact Lie group}
\author{Skip Garibaldi \and Robert M. Guralnick \and Eric M. Rains}
\maketitle
\setcounter{section}{17}
\section{Minimum values of characters: general case} \label{general.sec}

In the examples we have considered so far, $\min_{g \in G} \RE \chi(g)$ has been negative, which is easily seen to be a general phenomenon:

\begin{prop} \label{neg.val}
Let $\chi$ be the character of an irreducible (complex) representation of a
compact Lie group $G$.
\begin{enumeratea}
\item \label{neg.neg} If $\chi$ is nontrivial, then there exists $g\in G$ such that $\RE(\chi(g))<0$.
\item \label{neg.0} If $\chi(1) > 1$, then there exists $g \in G$ such that $\RE \chi(g) = 0$.
\end{enumeratea}
\end{prop}

\begin{proof}
Since $\chi$ is nontrivial, orthogonality implies that the
integral of $\chi$ relative to Haar measure is 0, and thus so is the
integral of $\RE\chi$. Since the integral of a nonzero nonnegative
continuous function is positive, \eqref{neg.neg} follows.

Now consider \eqref{neg.0}. The identity component $G^\ic$ is a normal subgroup of $G$, so its fixed subspace is an invariant submodule. It~follows that $G^\ic$ either acts trivially or without fixed points.
If $G^\ic$ acts trivially, then the action of $G$ factors through the finite group $G/G^\ic$, which is not cyclic (because $\chi$ is irreducible and $\chi(1) > 1$), so \eqref{neg.0} holds by Burnside \cite[p.\,115]{Burnside:tr}. If $G^\ic$ acts without fixed points, then the average value of $\RE \chi$ on $G^\ic$ is zero; since $\RE \chi(G^\ic)$ is connected and takes a positive value~$\chi(1)$, it~must also take a negative value, and we conclude \eqref{neg.0}.
\end{proof}

For self-dual representations, the statement of \ref{neg.val}\eqref{neg.0} says that there exists a $g \in G$ such that $\chi(g) = 0$, which
sounds similar to Theorem 1 in \cite{Se:trace}. Serre's statement adds the extra conclusion that $g$ can be chosen to be of finite order and does not require the self-dual hypothesis. See \cite{Se:zeros} for a proof of that result.

The aim of this section is to strengthen \ref{neg.val}\eqref{neg.neg} by proving the following uniform
bound on $\RE\chi / \chi(1)$:
\begin{thm} \label{neg.thm}
For any connected compact Lie group $G$, there is a constant $c_G>0$ such that for
any nontrivial, irreducible character $\chi$, there exists $g\in G$ such that
$\RE\chi(g) \le -c_G\chi(1)$.
\end{thm}

We treated the case $G = \SU(2)$ in Proposition \href{https://doi.org/10.5802/jep.310}{16.3 of the source article} by explicit calculation and found the optimal value of $c_{\SU(2)}$. (That calculation suggests that there is little hope in making the optimal $c_G$ explicit for general $G$.)
As we saw in that proof, the claim in the theorem
is closely related to the behavior of the characters themselves as the
weight becomes large, suggesting that we should consider limits of the form
\begin{equation} \label{X.lim}
\lim_{n\to\infty}
\frac{\chi_{\lambda^{(n)}}(e^{it/d_n})}{\chi_{\lambda^{(n)}}(1)}
\end{equation}
in which the sequence $\lambda^{(n)}/d_n$ converges, where $\lambda^{(n)}$ is a dominant weight and $d_n \in \R$ is positive.
When the limit of
$\lambda^{(n)}/d_n$ is in the interior of the fundamental chamber, this
limit is straightforward to compute, and leads us to define the function
\begin{equation} \label{Xdef}
X(s,t):=
\biggl( \prod_{r\in \Phi^+} \frac{r\cdot \rho}{(r\cdot s)(r\cdot t)} \biggr)
\biggl( \sum_{w\in W} \sgn(w) e^{s\cdot wt} \biggr)
\end{equation}
for $\Phi^+$ the positive roots and $s,t$ in the Cartan algebra $\mft(\C)$ of the \emph{complex} form of~$G$. By convention, we~take the
split form, so that the Lie algebra of $T$ is $i{\mft}(\R)$. In~the right side of \eqref{Xdef}, the products $r \cdot s$ and $r \cdot t$ are canonically determined by the root system. Since $G$ is compact, there is a $G$-invariant inner product on $\mfg$ \cite[\S{IX.1.3}, Prop.\,1]{Bou:g7}; pick one and use it to calculate $s \cdot wt$. The inner product defines an isomorphism of $\mft(\C)$ with its dual, and therefore also an inner product among elements of the dual which we use to evaluate $r \cdot \rho$. Up to scaling one of the arguments by $i$, $X$ is the Fourier transform of the Duistermaat-Heckman measure \cite{EtingofRains}.

Proposition \ref{167} below makes explicit the connection between \eqref{X.lim} and \eqref{Xdef}.
We will now take the analytic continuation of $X$,
which has the
advantage that it lets us extend the
function to the reflection hyperplanes.

\begin{prop} \label{Xprop}
The function $X(s,t)$ extends to an entire function of $\mft(\C)^2$
that satisfies the symmetries
\begin{enumeratea}
\item $X(s,t)=X(t,s)$,
\item $X(s,ut)=X(us,t)$ for $u\in \C$,
\item $X(s,wt)=X(s,t)$ for $w\in W$, and
\item $X(0,0)=1$.
\end{enumeratea}
\end{prop}

The proof will use the following lemma, which is implicit already in the development of the Weyl character formula. It~is a special case of a more general statement where~$\C^n$ is replaced by a connected complex manifold $X$ and the reflection $r$ is replaced by an automorphism of $X$ of finite order.

\begin{lem} \label{value}
Let $r\in \GL_n(\C)$ be a reflection, and let $f(z)$ be an $r$-invariant meromorphic function on an $r$-invariant region in $\C^n$
containing points fixed by $r$. Then the valuation of $f$ along the hyperplane $r\cdot z=0$ is even.
\end{lem}

\begin{proof}
This claim is clearly basis-independent, so we may assume that $r$ fixes the first $n-1$ coordinates and negates the last coordinate, so that we are claiming that $f$ has even order along the hyperplane $z_n=0$.
Let $z'$ be a fixed point of $r$ in the region where $f$ is defined, such that no other polar divisor of $f$ contains $z'$. Then $f$ has a Laurent series expansion in some polydisc around $z'$ of the form
\[
f(z) = \sum_{m_1,\dots,m_{n-1}\ge 0,m_n\ge v} c_{\vec{m}} \prod_{1\le i\le n-1} (z_i-z'_i)^{m_i} z_n^{m_n},
\]
where $v$ is the valuation, so that there is at least one nonzero term with $m_n=v$. But this must be invariant under negating $z_n$, so that any monomial in the expansion must have $m_n$ even, implying that $v$ is even as well.
\end{proof}

\begin{proof}[Proof of Proposition \ref{Xprop}]
The function $X(s,t)$ is certainly analytic on the complement of the
reflection hyperplanes $r\cdot s = 0$ and $r\cdot t = 0$, where it clearly
satisfies the symmetries $X(s,t)=X(t,s)$; $X(s,ut)=X(us,t)$ for $u\ne 0$;
and $X(s,wt)=X(s,t)$.

In particular, these same symmetries hold for the function $X(s,t)$, which
by inspection is a meromorphic function with at most a simple pole along
any reflection hyperplane. But the Lemma \ref{value} tells us that the order of such a pole must be even, and thus $X(s,t)$ actually has nonnegative valuation along any reflection hyperplane. Thus $X(s,t)$ is a meromorphic function with trivial polar divisor, so is holomorphic.

The symmetries continue to hold on the entire continuation, and thus in
particular we may take the limit $u\to 0$ in $X(s,ut)=X(us,t)$ to obtain
$X(s,0)=X(0,t)$ for all $s$ and $t$, and thus in particular both are equal to
$X(0,0)$.

It remains only to verify that $X(0,0)=1$. From the degenerate scaling
symmetry, this is equal to $X(\rho,0)$ for $\rho$ the half-sum of the positive roots, and thus to $\lim_{t\to 0}
X(\rho,t)$. We~moreover have the product formula
\[
X(\rho,t) =
\prod_{r\in \Phi^+} \frac{e^{r\cdot t/2}-e^{-r\cdot t/2}}{r\cdot t}
\]
in which each factor converges to $1$, so that $X(\rho,0)=1$ as required.
\end{proof}

\begin{rmk}
Note that $X$ only depends on the Weyl group, even though the root system appears in the definition of $X$.
Rescaling the roots (including
swapping short and long roots in non-simply-laced cases) only
multiplies $X$ by a constant, which must be 1 by virtue of $X(0,0)=1$.
\end{rmk}

Since $X(s,t)$ differs by simple factors from the Weyl numerator, we~can
express the normalized characters in terms of $X(s,t)$.

\begin{lem} \label{X.dom}
For any $t\in \mft(\R)$ and any dominant weight $\lambda$, one has
\[
\frac{\chi_{\lambda}(e^{it})}{\chi_{\lambda}(1)}
=
\frac{X(\lambda+\rho,it)}{X(\rho,it)}.
\]
\end{lem}

\begin{proof}
Writing out $X(\lambda+\rho, it)/X(\rho, it)$ using the definition \eqref{Xdef}, gives $\chi_\lambda(e^{it})$ divided by
\[
\prod_{r\in \Phi^+} \frac{r\cdot (\lambda+\rho)}{r\cdot \rho} = \chi_\la(1),
\]
where the equality is the Weyl dimension formula.
\end{proof}

This lets us show that $X(s,t)$ is indeed the limiting character even on
the boundary of the fundamental chamber.

\begin{prop} \label{167}
Let $(\lambda^{(n)},d_n)$ be a sequence of pairs where each $\lambda^{(n)}$ is a dominant weight and $d_n \in \R$ is such that
$\lambda^{(n)}/d_n$ converges and $d_n \to \infty$. Then
\[
\lim_{n\to\infty}
\frac{\chi_{\lambda^{(n)}}(e^{it/d_n})}{\chi_{\lambda^{(n)}}(1)}
=
X\bigl(\lim_{n\to\infty} \lambda^{(n)}/d_n,it \bigr)
\]
for all $t\in \mft(\R)$.
\end{prop}

\begin{proof}
From Lemma \ref{X.dom}, one has
\begin{align*}
\lim_{n\to\infty}
\frac{\chi_{\lambda^{(n)}}(e^{it/d_n})}{\chi_{\lambda^{(n)}}(1)}
& =
\lim_{n\to\infty}
\frac{X(\lambda^{(n)}+\rho,it/d_n)}{X(\rho,it/d_n)} \\
&=
\lim_{n\to\infty}
\frac{X(\lambda^{(n)}/d_n+\rho/d_n,it)}{X(\rho/d_n,it)},
\end{align*}
with the second equality following from the scaling symmetry of $X$.
Since $X$ is continuous, we~can pull the limit inside, at which point
the claim follows by noting that $X(0,it)=1$.
\end{proof}

To establish the desired bound on $\min\RE\chi/\chi(1)$, we~will
show after some work that $X(s,it)$ takes negative values for any nonzero $s\in \mft(\R)$ (Lemma \ref{t.neg}). We~write~$\lambda$ for the Haar measure on $\mft(\R)$, since it is essentially just Lebesgue measure.
To~avoid issues with the normalization of the Haar measure, the root
system, and the $G$\nobreakdash-invariant inner product, define a scaling constant
\[
Z := \int_{t\in \mft(\R)} \biggl( \prod_{r\in \Phi^+}(r\cdot t)^2 \biggr) e^{-t\cdot
t/2}
\lambda(dt).
\]

We also use the following mild restatement of the well-known formula for the Fourier or Laplace transform of a Gaussian:
\begin{lem} \label{fourier}
For any $s\in \mft(\C)$, one has
\[
\int_{t \in \mft(\R)} e^{s\cdot t} e^{-t\cdot t/2} \lambda(dt)
=
e^{s\cdot s/2} \int_{t \in \mft(\R)} e^{-t\cdot t/2} \lambda(dt).
\]
\end{lem}

\begin{proof}
The integral on the left is absolutely convergent for any $s\in \mft(\C)$, and uniformly so given any bound on $\left|\RE(s)\right|$, so describes an entire function of $s$. Since the right-hand side is manifestly entire, it~suffices to verify that the two sides agree when $s$ is real. But for $s\in \mft(\R)$, we~may complete the square on the left by translating~$t$ by $s$ to obtain
\begin{align*}
\int_{t\in \mft(\R)} e^{s\cdot t} e^{-t\cdot t/2} \lambda(dt)
&=
\int_{t\in -s+\mft(\R)} e^{s\cdot s/2} e^{-t\cdot t/2} \lambda(dt)\\
&=
e^{s\cdot s/2} \int_{t\in \mft(\R)} e^{-t\cdot t/2} \lambda(dt),
\end{align*}
as required.
\end{proof}

\begin{prop} \label{158}
For $s,s'\in \mft(\C)$, one has
\[
\frac{1}{Z}
\int_{t\in \mft(\R)} X(s,it)X(s',it)
\biggl( \prod_{r\in \Phi^+} (r\cdot t)^2 \biggr)
e^{-t\cdot t/2}
\lambda(dt)
= 
e^{-s\cdot s/2-s'\cdot s'/2}
X(s,s').
\]
\end{prop}

\begin{proof}
Since the right-hand side is analytic in $s$, $s'$, it~suffices to show
the equality when neither of $s$ or $s'$ lies in any reflection
hyperplane. We~can then write
\begin{multline*}
\int_{t\in \mft(\R)} X(s,it)X(s',it)
\biggl( \prod_{r\in \Phi^+} (r\cdot t)^2\biggr)
e^{-t\cdot t/2}
\lambda(dt)
\\
=\biggl(\prod_{r\in \Phi^+} \frac{(r\cdot \rho)^2}{(r\cdot s)(r\cdot s')}\biggr)
\int_{t\in \mft(\R)}
\sum_{w,w'\in W}
\sgn(ww')
e^{i(ws+w's')\cdot t}
e^{-t\cdot t/2}
\lambda(dt).
\end{multline*}
Applying Lemma \ref{fourier} gives:
\begin{multline*}
\int_{t\in \mft(\R)}
\sum_{w,w'\in W}
\sgn(ww')
e^{i(ws+w's')\cdot t}
e^{-t\cdot t/2}
\lambda(dt)\\[-5pt]
=\sum_{w,w'\in W}
\sgn(ww')
e^{-(ws+w's')\cdot(ws+w's')/2}
\int_{t\in \mft(\R)}
e^{-t\cdot t/2}
\lambda(dt).
\end{multline*}
We then observe that
\begin{multline*}
\sum_{w,w'\in W}
\sgn(ww')
e^{-(ws+w's')\cdot (ws+w's')/2}
\\[-5pt]
= e^{-s\cdot s/2}e^{-s'\cdot s'/2}
\sum_{w,w'\in W} \sgn(ww') e^{ws\cdot w's'}
\\
= |W| \,
e^{-s\cdot s/2} e^{-s'\cdot s'/2}
\biggl( \prod_{r\in \Phi^+}\frac{(r\cdot s)(r\cdot s')}{r\cdot \rho} \biggr)
X(s,s').
\end{multline*}
We thus conclude that
\[
\int_{t\in \mft(\R)} X(s,it)X(s',it)
\biggl( \prod_{r\in \Phi^+} (r\cdot t)^2 \biggr)
e^{-t\cdot t/2}
\lambda(dt)
\]
is proportional to
\[
e^{-s\cdot s/2}e^{-s'\cdot s'/2}
X(s,s'),
\]
and can recover the constant by setting $s=s'=0$.
\end{proof}

The proposition is already enough to give us a weak form of orthogonality sufficient
to prove the desired negativity result.

\begin{cor} \label{limit.0}
If $\RE(s\cdot s)>0$, then
\[
\lim_{u\to\infty}
\int_{t\in \mft(\R)} X(s,it)
\biggl( \prod_{r\in \Phi^+} (r\cdot t)^2 \biggr)
e^{-t\cdot t/2u}
\lambda(dt)
=
0.
\]
\end{cor}

\begin{proof}
Setting $s'=0$ and abbreviating $p(t) := \prod_{r\in \Phi^+} (r\cdot t)^2$ in Proposition \ref{158} gives
\[
\frac{1}{Z}
\int_{t\in \mft(\R)} X(s,it)
\,p(t)\,
e^{-t\cdot t/2}
\lambda(dt)
=
e^{-s\cdot s/2},
\]
and thus
\[
\frac{1}{Z}
\int_{t\in \mft(\R)} X(u^{1/2}s,it)
\, p(t)\,
e^{-t\cdot t/2}
\lambda(dt)
=
e^{-u(s\cdot s)/2}.
\]
We can use the scaling symmetry of $X$ to move the $u^{1/2}$ to the other
argument and then do a change of variables to move it to the other parts
of the integrand. We~thus find
\begin{multline*}
\frac{1}{Z}
\int_{t\in \mft(\R)} X(u^{1/2}s,it)
\,p(t)\,
e^{-t\cdot t/2}
\lambda(dt)\\
=u^{-|\Phi^+|/2-\dim\mft/2}
\frac{1}{Z}
\int_{t\in \mft(\R)} X(s,it)
\, p(t) \,
e^{-t\cdot t/2u}
\lambda(dt),
\end{multline*}
and thus
\[
\frac{1}{Z}
\int_{t\in \mft(\R)} X(s,it)
\, p(t) \,
e^{-t\cdot t/2u}
\lambda(dt)
=
u^{(|\Phi^+|+\dim\mft)/2}
e^{-u(s\cdot s)/2},
\]
which converges to $0$ as required.
\end{proof}

\begin{lem} \label{t.neg}
For any nonzero $s\in \mft(\R)$, there exists $t\in \mft(\R)$
such that $\RE X(s,it)<0$.
\end{lem}

\begin{proof}
Suppose otherwise, and consider the integral
\[
\frac{1}{Z}
\int_{t\in \mft(\R)}
\RE X(s,it) \,
\biggl(\prod_{r\in \Phi^+} (r\cdot t)^2 \biggr)
e^{-t\cdot t/2u}
\lambda(dt)
\]
as a function of $u$. On the one hand, the integrand is nonnegative,
not identically~0, and monotonically increasing in $u$, and thus the
integral itself must be a positive, monotonically increasing, function of
$u$. On the other hand, Corollary \ref{limit.0} tells us that the
integral converges to 0, giving a contradiction.
\end{proof}

\begin{proof}[Proof of Theorem \ref{neg.thm}]
If $G$ has positive-dimensional center $Z$, then for any character $\chi$
which is nontrivial on $Z^\ic$, there is an element $z\in Z^\ic$ such that
$\chi(z)=-\chi(1)$. Since one has the universal bound $\RE(\chi(g))\ge
-\chi(1)$, it~follows that any constant that works for $G/Z^\ic$ will work
for $G$ as well. We~thus reduce to the case that $G$ is semisimple.

In the semisimple case, replacing $G$ by its universal cover only makes the conclusion harder to satisfy, and thus we may assume $G$ simply connected. Moreover, if $G=G_1\times G_2$, then any irreducible character of $G$ factors as $\chi_1\boxtimes \chi_2$; it follows that if the theorem holds for both factors with constants $c_1$, $c_2$, then it holds for $G$ with constant at least $\min(c_1,c_2)$. (Take $g$ to be the identity on one of the two factors.) We thus reduce to the case that $G$ is simply connected and simple.

Now, if the claim fails for $G$, then there exists a sequence $\lambda^{(n)}$
of dominant weights such that\vspace*{-3pt}
\[
\limsup_{n\to\infty} \inf_{g\in G}
\frac{\RE\chi_{\lambda^{(n)}}(g)}{\chi_{\lambda^{(n)}}(1)}\ge 0.
\]
We may replace this by a subsequence to replace the limsup by a limit.
Moreover, since spheres are compact, we~may pass to a further subsequence
to ensure that the sequence $\lambda^{(n)}/|\lambda^{(n)}|$ converges,
say to $s$. We~then have by Proposition \ref{167}, for any $t\in \mft(\R)$\vspace*{-3pt}
\[
\lim_{n\to\infty}
\frac{\chi_{\lambda^{(n)}}(e^{it/|\lambda^{(n)}|})}{\chi_{\lambda^{(n)}}(1)}
=
X(s,t),
\]
and by Lemma \ref{t.neg} may choose $t$ so that $\RE X(s,t)<0$. But this gives a
contradiction:\vspace*{-3pt}
\[
\inf_{g\in G}
\frac{\RE\chi_{\lambda^{(n)}}(g)}{\chi_{\lambda^{(n)}}(1)}
\le
\frac{\RE \chi_{\lambda^{(n)}}(e^{it/|\lambda^{(n)}|})}{\chi_{\lambda^{(n)}}(1)}
\]
but the right-hand side has negative limit while the left-hand side has
nonnegative limit by assumption.
\end{proof}
\begin{thebibliography}{BGR23}
\bibitem[Bur03]{Burnside:tr} W. Burnside – “On an arithmetical theorem connected with roots of unity and its application to group characteristics”, Proc. London Math. Soc. 1 (1903), p. 112–116.
\bibitem[Ser04]{Se:trace} J.-P. Serre, “On the values of the characters of compact Lie groups”, in Oberwolfach Reports, vol. 1, EMS Publishing House, Zürich, 2004, p. 666–667.
\bibitem[Ser25]{Se:zeros} J.-P. Serre, “Zéros de caractères”, Enseign. Math. (2) 71 (2025), p. 433–457.
\bibitem[Bou05]{Bou:g7} N. Bourbaki, Lie groups and Lie algebras: Chapters 7–9, Springer-Verlag, Berlin, 2005.
\bibitem[ER25]{EtingofRains} P. Etingof \& E. Rains – “Bounds for the asymptotic characters of simple Lie groups”, Indag. Math. (2025), online, arXiv:2405.10341.
\end{thebibliography}
\end{document}
